%! AP Statistics — Unit 5, Lesson 5.2 — Student Packet Cover Sheet
\documentclass[10pt]{article}
\usepackage{apstats-article}
\usepackage{apstats-boxes}
\usepackage{ltablex}
\keepXColumns

\begin{document}

% Full-bleed navy banner
\begin{tikzpicture}[remember picture, overlay]
  \fill[navy] (current page.north west)
    rectangle ([yshift=-0.9in]current page.north east);
\end{tikzpicture}
\vspace{-0.8in}

\begin{center}
  {\color{white}\textbf{\LARGE AP Statistics}}\\[4pt]
  {\color{skymid}\large\textbf{Unit 5: Sampling Distributions}}\\[6pt]
  {\color{white}\textbf{\large Lesson 5.2 \quad The Normal Distribution, Revisited}}
\end{center}

\vspace{0.22in}
\namedateperiod
\vspace{0.18in}

\begin{learningtargetbox}
\small
\begin{enumerate}[leftmargin=1.5em, itemsep=5pt]
  \item I can calculate the probability that a continuous random variable with a normal
        distribution falls in a given interval using z-scores and/or technology.
        \textit{(VAR-6.A)}
  \item I can determine the boundary value(s) of an interval in a normal distribution
        associated with a given area. \textit{(VAR-6.B)}
  \item I can assess whether the normal distribution is an appropriate model for a
        given distribution. \textit{(VAR-6.C)}
\end{enumerate}
\end{learningtargetbox}

\vspace{0.14in}

\begin{tocbox}
\small
\renewcommand{\arraystretch}{1.5}
\begin{tabularx}{\linewidth}{@{} c l X r @{}}
\rowcolor{sky}
\textbf{\#} & \textbf{Component} & \textbf{Description} & \textbf{Score} \\
\midrule
1 & Warm-Up        & z-scores and the 68--95--99.7 rule                      & \blank{1.2cm} \\
2 & Group Activity & Normal probability calculations across three contexts    & \blank{1.2cm} \\
3 & Exit Ticket    & Forward and inverse normal problems                      & \blank{1.2cm} \\
4 & Homework       & Normal probability practice + model assessment           & \blank{1.2cm} \\
\midrule
  & \multicolumn{2}{r}{\textbf{Total}} & \blank{1.2cm} \\
\end{tabularx}
\end{tocbox}

\end{document}
